\section{The SOLE Encoding}
\label{sec:sole}

\subsection{Problem Setting}

Let a machine block contain \(b\) bits, and define
\[
    B=2^b,
    \qquad
    [B]=\{0,1,\ldots,B-1\}.
\]
An ordinary block is an element of \([B]\). To make a stream
self-delimiting, SOLE adds a distinguished symbol \(\mathsf{eof}\), identified
with the value \(B\); the enlarged alphabet is therefore \([B+1]\).

Basic SOLE considers an odd number \(n\) of input blocks and assumes
\[
    b\geq 2\log_2 n+1,
    \qquad
    B\geq 2n^2.
\]
Under these assumptions it encodes the \(n\) data blocks, together with the
termination information, as \(n+2\) ordinary blocks. Its name abbreviates
\emph{Short-Odd Long-Even}: the first pass deliberately creates alternating
short and long ranges, which the second pass recombines.

\subsection{The First Pass: Short-Odd Long-Even}

For \(i=0,1,\ldots\), the first pass processes symbols in positions
\(2i+1\) and \(2i+2\). If their values are \(a_0,a_1\in[B+1]\), it combines
them as
\[
    z=a_1(B+1)+a_0.
\]
It then divides \(z\) by \(B-4i\), writing
\[
    z=q(B-4i)+r.
\]
The output pair is \((r,q)\), and the range conversion is
\[
    [B+1]^2
    \longrightarrow
    [B-4i]\times[B+4i+4].
\]
The parameter assumptions ensure the capacity inequality
\[
    (B+1)^2
    \leq
    (B-4i)(B+4i+4).
\]
Consequently \(r\in[B-4i]\) and \(q\in[B+4i+4]\). Quotient and
remainder uniquely determine \(z\), and division by \(B+1\) then recovers
\(a_0\) and \(a_1\), so the transformation is reversible.

\subsection{Regrouping and the Second Pass}

After the first pass the ranges alternate between values slightly below and
slightly above \(B\). The second pass shifts the pairing by one position, so
that a long value is coupled with the following short value. A typical pair
has the conversion
\[
    [B+4i]\times[B-4i]
    \longrightarrow
    [B]^2.
\]
Indeed,
\[
    (B+4i)(B-4i)=B^2-16i^2<B^2.
\]
Combining the two nonstandard values into one mixed-radix integer and taking
its quotient and remainder modulo \(B\) therefore produces two ordinary
blocks. The short and long ranges need not fit binary memory separately;
their product fits once they are regrouped.

\subsection{Termination and Prefix-Freeness}

Because \(n\) is odd, the appended \(\mathsf{eof}\) symbol occupies the even
position \(n+1\) and is processed with the final data block in the first pass.
If this last pair has zero-based pair index \(i\), the algorithm additionally
appends the value \(0\), regarded as an element of
\[
    [B-4i-4].
\]
This artificial value supplies the short range needed to pair with the final
long range in the second pass.

During decoding, the inverse local transformations recover enlarged-alphabet
symbols in order. The decoder stops immediately upon recovering
\(\mathsf{eof}\), so the artificial zero is not part of the original message.
Property~3 of SOLE bounds the required decoding window by a constant number of
nearby output blocks. In particular, the window that reveals
\(\mathsf{eof}\) lies within the available encoding, so the decoder never has
to read past its end. The resulting code is therefore self-delimiting and
prefix-free.

\subsection{A Concrete Example}

Take
\[
    n=3,
    \qquad
    B=32,
    \qquad
    (a_1,a_2,a_3)=(5,7,31).
\]
Here \(\mathsf{eof}=32\), and the augmented sequence is
\((5,7,31,32,0)\). For the first pair, Pass~1 gives
\[
    7\cdot 33+5=236=7\cdot 32+12,
\]
and hence \((12,7)\). For the second pair,
\[
    32\cdot 33+31=1087=38\cdot 28+23,
\]
and hence \((23,38)\). The intermediate sequence is therefore
\[
    (12,7,23,38,0).
\]

Pass~2 leaves the first \(12\) in place and regroups the remaining values.
For \((7,23)\), whose ranges have sizes \(36\) and \(28\),
\[
    23\cdot 36+7=835=26\cdot 32+3,
\]
which produces \((3,26)\). The final group gives
\[
    0\cdot 40+38=38=1\cdot 32+6,
\]
which produces \((6,1)\). Thus the complete output is
\[
    (12,3,26,6,1).
\]

\subsection{Locality and Updates}

Property~3 makes the locality quantitative. Output blocks \(2i\) and
\(2i+1\) are determined by input blocks \(2i-1,\ldots,2i+2\), while input
blocks \(2i+1\) and \(2i+2\) can be recovered from output blocks
\(2i,\ldots,2i+3\). Thus a constant number of input blocks determines a
constant number of output blocks, and conversely.

Changing one input block consequently requires recomputing only a constant
neighborhood of the encoding. Appending a block or truncating the stream can
likewise be reduced to a constant number of local writes near the end. This
worst-case locality distinguishes SOLE from a global mixed-radix conversion.

\subsection{Variants and Practical Remarks}

The basic presentation favors a transparent termination argument and uses
\(n+2\) output blocks. A tighter termination treatment reduces the boundary
waste and yields an \(n+1\)-block version. The paper also gives the symmetric
LOSE encoding, which reverses the short/long placement when that orientation
is more convenient.

For practical block sizes such as \(b=32\) or \(b=64\), all operands fit in a
constant number of machine words, and the quotient--remainder transformations
can be implemented with standard word arithmetic. SOLE is therefore not
merely an asymptotic counting argument.

\subsection{From SOLE to Information Carriers}

The non-power-of-two values temporarily left in SOLE's short and long ranges
are instances of a \emph{spill}: information that has not yet been committed
to standard binary memory. Later blocks help commit an old spill while
leaving a new, smaller spill. Section~\ref{sec:information-carriers} abstracts
this mechanism into the information-carrier lemma, independent of SOLE's
particular alternating ranges.
